%%% Local Variables:
%%% mode: latex
%%% TeX-master: "../main"
%%% End:

\chapter[多任务强化学习总体框架与混合地形统一环境]{多任务强化学习总体框架与混合地形统一环境}
\label{cha:framework}

本章针对两轮足机器人"多地形、多步态、多算法"的综合任务需求，提出"分训-整合"两阶段的多任务强化学习总体框架：第一阶段分别训练7个步态专家策略，第二阶段以学习式统一调度器将各专家整合至混合地形统一环境中，实现按地形自适应切换步态。本章首先给出总体设计思路，然后阐述混合地形统一环境设计、步态专家库划分与学习式调度器的总体方案，最后说明三大核心问题（灾难性遗忘、奖励冲突、局部最优）的应对策略。

\section{总体设计思路}

\subsection{设计需求分析}

综合前文分析，总体框架需要满足以下设计需求：

（1）\textbf{多步态覆盖}：支持滚动行走、侧向移动、粗糙地形行走、抬腿行走、跳跃、后空翻、抬腿上台阶与单腿平衡等步态能力；

（2）\textbf{多地形适应}：覆盖平面、粗糙、碎石、斜坡、台阶、弧度、独木桥、木桩与狭窄通道等多种地形，并能在混合地形环境中自动切换步态；

（3）\textbf{多算法融合}：支持单PPO、PPO+SLIP、SLIP+PPO+VMC、N3PO等算法范式的统一接入与对比；

（4）\textbf{多任务整合}：系统性地缓解灾难性遗忘、奖励冲突与局部最优三大问题。

\subsection{"分训-整合"两阶段范式}

针对单一策略难以同时掌握多种步态、直接统一训练又易遗忘与冲突的难题，本文提出"分训-整合"两阶段范式，如图~\ref{fig:paradigm}所示。

\textbf{第一阶段（分训）}：将综合任务分解为7类典型步态子任务，每类步态在对应地形的专用环境中独立训练，得到7个参数解耦的步态专家策略网络。各专家互不干扰，训练目标单一明确，从根本上规避了多任务共同训练导致的灾难性遗忘与奖励冲突。

\textbf{第二阶段（整合）}：将7个步态专家作为底层选项集合，设计学习式统一调度器。调度器在混合地形统一环境中训练，学习依据地形感知与本体状态自动选择或混合激活合适的步态专家，实现机器人从起点到终点的完整遍历与按地形自适应步态切换。

\begin{figure}[!htbp]
    \centering
    \includegraphics[width=0.95\textwidth]{WUST-flash.pdf}
    \bicaption{"分训-整合"两阶段范式示意图（替换为自绘架构图）}{Schematic of the two-phase "train-then-integrate" paradigm (replace with your own figure)}
    \label{fig:paradigm}
\end{figure}

\section{混合地形统一仿真环境设计}

\subsection{地形库构成}

为验证多步态整合与按地形自动切换效果，本文构建包含多种地形单元的混合地形统一仿真环境，地形库如表~\ref{tab:terrain-lib}所示。

\begin{table}[!htbp]
    \centering
    \bicaption{混合地形统一环境的地形库构成}{Terrain library of the unified mixed-terrain environment}
    \resizebox{0.95\textwidth}{!}{
        \begin{tabular}{clp{5.6cm}}
            \toprule[1.5pt]
            地形单元 & 关键特征 & 设计目的 \\
            \midrule[0.75pt]
            平面地形 & 平整连续 & 滚动行走、侧向移动基础验证 \\
            粗糙/碎石地形 & 高度起伏、接触随机 & 验证抬腿点足与粗糙行走专家 \\
            斜坡地形 & 坡角变化 & 验证地形自适应腿长调节 \\
            台阶地形（低/高） & 不同台阶高度 & 低阶抬腿跨越、高阶跳跃上去 \\
            独木桥地形 & 窄支撑面 & 验证单腿平衡专家 \\
            木桩/间隙地形 & 离散支撑、有间隙 & 验证跳跃专家跨越间隙 \\
            狭窄通道 & 转向空间受限 & 验证侧向移动专家 \\
            \bottomrule[1.5pt]
        \end{tabular}
    }
    \label{tab:terrain-lib}
\end{table}

\subsection{地形-步态对应关系}

按地形自动切换步态是整合阶段的核心目标，本文明确各典型地形与最适步态专家的对应关系，作为调度器训练的期望行为与验收基准，如表~\ref{tab:terrain-gait}所示。

\begin{table}[!htbp]
    \centering
    \bicaption{典型地形与最适步态的对应关系}{Correspondence between typical terrains and preferred gaits}
    \resizebox{0.95\textwidth}{!}{
        \begin{tabular}{lcp{4.6cm}}
            \toprule[1.5pt]
            地形场景 & 采用步态 & 理由 \\
            \midrule[0.75pt]
            粗糙/碎石地形 & 抬腿点足行走 & 减少轮端颠簸、提高越障能力 \\
            低台阶地形 & 抬腿跨越 & 小位移精细跨越，稳定高效 \\
            高台阶地形 & 跳跃上去 & 超出抬腿范围，需要弹道起跳 \\
            狭窄通道地形 & 侧向移动 & 横向位移，避免大半径转向 \\
            独木桥地形 & 单腿平衡 & 缩小有效支撑宽度，精确过桥 \\
            木桩/间隙地形 & 跳跃通过 & 跨越支撑间隙 \\
            平坦开阔地形 & 滚动行走 & 高效平稳 \\
            \bottomrule[1.5pt]
        \end{tabular}
    }
    \label{tab:terrain-gait}
\end{table}

\subsection{领域随机化}

为提升策略与调度器的鲁棒性、抑制局部最优，统一环境在训练中随机化机器人质量、质心位置、摩擦系数、地形参数（坡角、台阶高度、间隙宽度）等物理量，并施加随机外部扰动力，迫使策略与调度器学习不依赖特定参数的鲁棒控制与决策方式\cite{tobin2017,peng2018}。

\section{步态专家库的独立训练}

\subsection{7个步态专家的划分}

依据综合任务需求，本文将步态能力划分为7个步态专家，如表~\ref{tab:experts}所示。每个专家对应一个独立的策略网络，在对应地形的专用环境中训练。

\begin{table}[!htbp]
    \centering
    \bicaption{步态专家库划分与训练配置}{Gait expert library and training configuration}
    \resizebox{0.95\textwidth}{!}{
        \begin{tabular}{clll}
            \toprule[1.5pt]
            序号 & 步态专家 & 适用地形 & 训练算法 \\
            \midrule[0.75pt]
            1 & 滚动行走专家（含侧向移动） & 平面、狭窄通道 & PPO \\
            2 & 粗糙地形行走专家 & 粗糙、碎石、斜坡 & PPO \\
            3 & 抬腿行走专家 & 粗糙、木桩 & PPO \\
            4 & 跳跃专家 & 高台阶、间隙 & PPO / PPO+SLIP / SLIP+PPO+VMC \\
            5 & 后空翻专家 & 平面（特殊动作） & PPO（确定性翻转） \\
            6 & 抬腿上台阶专家 & 台阶（低） & N3PO \\
            7 & 单腿平衡专家 & 独木桥、狭窄 & PPO \\
            \bottomrule[1.5pt]
        \end{tabular}
    }
    \label{tab:experts}
\end{table}

\subsection{专家独立训练与验收}

各步态专家采用统一的观测组织方式（本体状态 + 地形感知 + 任务指令）与训练流程，在专用环境中训练至各自验收标准（如站立时长、速度跟踪误差、成功率等）。由于各专家参数完全解耦，独立训练过程互不影响，从结构上规避了多步态连续学习带来的灾难性遗忘。第4-6章将分别阐述行走类、精细平衡类与高动态敏捷类专家的具体训练细节。

\subsection{观测、动作与执行频率设计}

观测空间按"本体状态 + 地形感知 + 任务指令"三部分组织，各专家依据任务特点动态扩展维度（如跳跃专家追加起跳标志与飞行阶段信息，单腿平衡专家追加FSM状态与轮地接触历史）。动作空间定义为各关节目标角与轮端目标速度（或力矩）。基础行走类专家采用100Hz控制率，高动态专家（跳跃、后空翻）采用200Hz控制率，热启动迁移时保持控制率对齐。

\section{学习式统一调度器总体方案}

\subsection{调度器任务建模}

统一调度器将"按地形选择步态"建模为分层强化学习问题：底层为7个步态专家构成的选项集合，高层为调度策略$\pi_h$。调度器在时刻$t$接收地形感知特征$\boldsymbol{\phi}_t$与本体状态$\boldsymbol{s}_t$，输出对专家$i$的选择概率
\begin{equation}
  \pi_h(i\,|\,\boldsymbol{\phi}_t,\boldsymbol{s}_t) = \frac{\exp\left(f_i(\boldsymbol{\phi}_t,\boldsymbol{s}_t)/\tau\right)}{\sum_{j=1}^{M}\exp\left(f_j(\boldsymbol{\phi}_t,\boldsymbol{s}_t)/\tau\right)},
  \label{eq:scheduler}
\end{equation}
式中，$f_i(\cdot)$为调度网络对专家$i$的得分，$M$为专家数量，$\tau$为温度系数。调度器可采用硬门控（激活得分最高的专家）或软门控（按概率加权混合专家动作）。

\subsection{调度器训练}

调度器在混合地形统一环境中训练，其高层奖励综合任务完成度、遍历速度与稳定性等指标。训练采用课程式场景编排：由单地形单元逐步过渡到多地形交替，最终在完整混合地形环境中训练，使调度器逐步学会"地形感知-步态选择"的映射关系。调度器训练过程中，各步态专家可保持冻结以验证调度本身的有效性，也可在后期进行联合微调以进一步提升整体性能。

\subsection{与三大问题的关系}

"分训-整合"范式与三大核心问题的关系如下：各专家参数解耦避免了灾难性遗忘；各专家独立奖励与调度器高层目标分离缓解了奖励冲突；混合地形统一环境的细粒度评估与领域随机化抑制了局部最优。具体机制在第3.5节详细阐述。

\section{多任务整合的三大核心问题应对}

\subsection{灾难性遗忘：参数解耦与分训}

传统的单一网络按序学习多步态时，新步态的参数更新会覆盖旧步态知识，导致灾难性遗忘\cite{mccloskey1989,kirkpatrick2017}。本文通过"分训-整合"范式从结构上化解该问题：7个步态专家各自拥有独立的网络参数，分训阶段互不干扰；整合阶段调度器仅在专家之间进行选择，不直接修改专家参数（除非联合微调），因此任何步态能力的习得都不会破坏其他步态能力。

\subsection{奖励冲突：专家奖励与调度器目标分离}

多步态任务中"跳得高""走得稳""姿态受限"等子目标相互矛盾，若合并为单一奖励函数将难以优化。本文将该冲突通过分层结构分离：每个专家在各自的专用环境中拥有专属的、内部自洽的奖励函数（如跳跃专家专注跳跃高度与落地稳定）；调度器拥有独立的高层目标（如遍历完成度），不直接叠加专家内部奖励。通过"底层各管各的目标、高层只管选谁干"，奖励冲突被结构性消解。

\subsection{局部最优：统一评估与领域随机化}

为规避策略"抖动取巧"等局部最优，本文一方面建立统一的细粒度评估指标体系（平移/航向跟踪误差、成功率、轮端净空、落地冲击、关节力矩饱和度等），从多维度揭示策略行为是否走捷径；另一方面通过强领域随机化迫使策略学习鲁棒控制而非依赖特定参数的模式。

\section{多维评价与分析体系}

本文建立贯穿分训与整合两个阶段的统一评价体系：分训阶段评估各专家的步态性能与算法对比指标；整合阶段评估调度器在混合地形环境中的遍历成功率、步态切换次数、切换过程姿态扰动以及各地形段落的通过情况，并以多维度雷达图与统计表形式呈现。三大问题的消融验证（去调度直接统一训练、去领域随机化等）也基于该评价体系开展。

\section{本章小结}

本章提出"分训-整合"两阶段的多任务强化学习总体框架，设计了混合地形统一仿真环境与地形-步态对应关系，划分了7个步态专家并阐述其独立训练方案，给出了学习式统一调度器的总体建模与训练思路，并从参数解耦、奖励分离、统一评估三个层面阐述了三大核心问题的应对策略。后续第4-6章阐述步态专家训练细节，第7章阐述调度器设计，第8章进行统一系统集成与整体验证。
